% Chapter 0 -- Abstract
% Self-contained section file: no preamble, no \documentclass.
% Mirrors the template's bilingual dual-abstract slot (Thai p.5 + English p.6 + keyword line);
% English written here per the project's single-language directive.
% HONESTY RULE: this is a methodology / research-proposal abstract. No experiments have been
% run and no result numbers of ours are reported. Literature anchors are attributed as background.

\chapter*{Abstract}
\addcontentsline{toc}{chapter}{Abstract}
\label{ch:abstract}

This is a methodology and research proposal: no experiments have yet been
conducted, and the document reports no empirical results of our own. We
propose a \emph{trustworthiness-first} feline pain scorer that is deliberately
\emph{not} another cat-pain detector. Automated Feline Grimace Scale (FGS) work
has established strong in-domain detection numbers --- the FGS itself validates
at the cited reference level of AUC~$0.94$ with sensitivity~$90.7\%$ and
specificity~$86.6\%$ above a $>0.39$ threshold~\cite{evangelista2019}, with
subsequent landmark-and-geometry pipelines reporting comparable
accuracy~\cite{steagall2023, feighelstein2023} --- so a further accuracy
re-implementation would contribute nothing. Those figures are background only,
attributed to their authors, and are never restated as ours. Instead, the
\emph{headline} contribution is a single portable method, intended to transfer
to the next facial-pain corpus rather than to characterize a single dataset: a
\emph{power-aware, one-directional capture-condition confound-attribution
protocol} (two legs: an FGS-versus-background gap statistic, FGS-BG-Gap, together with a
per-AU energy-based pointing-game measure, EBPG read as per-AU
saliency-as-confound-evidence; a fully-coded VLM judge-bias axis is exercised on
synthetic controls only and named as future work, not part of the headline) that any future
corpus can re-run as a one-directional audit, reporting ``no confound detected
at this power'' rather than asserting the absence of confounding~\cite{adebayo2022posthoc}. This protocol
is validated today on planted positive and negative controls. Supporting it,
and explicitly subordinate to it, is a \emph{guarded} \emph{VLM-as-AU-rater}
reliability check: it scores a frozen vision--language model's per-action-unit
(per-AU) weak labels against a veterinary anchor by quadratic-weighted Cohen's
kappa, gated on the \emph{lower bound} of the kappa confidence interval rather
than its point estimate. We frame this check as inspected-not-validated rather
than as a co-equal second contribution: its result is pending an independent
per-AU veterinary anchor that does not yet exist (the current binary
pain/no-pain labels cannot yield $0/1/2$ AU ground truth), so it ships today as
a runnable reliability check rather than a reported number, and is retained as
kill-tree insurance --- the sole-survivor headline should the confound leg
degrade. The kappa CI-lower-bound gate is textbook clinimetrics and the
rubric-paraphrase independence guard is published practice; neither is claimed
as a novel increment. The novelty we do claim is the \emph{assembly}, the
per-AU EBPG-as-confound-evidence reading, and the instantiation of
equivalence-style audit hygiene at a pre-registered power --- never any
individual primitive nor the underlying power and equivalence statistics
themselves, which we credit to their
sources~\cite{huanghooker2026fairness,singh2023samplesize}.

The headline confound-attribution protocol is released as a dataset-agnostic,
importable, gate-enforced executable embodiment (the portable protocol surface);
the subordinate, guarded reliability check ships alongside it as an importable
module rather than as a co-equal second method: both are surfaced in
\texttt{src/protocols/} (re-exports plus thin adapters for
AU override / column mapping / a generic non-FGS ordinal mode with no $0.39$
assumption; see \texttt{src/protocols/adapters.py} and
\texttt{src/protocols/\_\_init\_\_.py}), exercised by the deletion-safe
zero-torch standalone test corpus
(\texttt{src/protocols/standalone\_test\_corpus.py}, which exercises the
confound protocol end-to-end on synthetic controls, with the planted
above-threshold positive-recovery assertions held in
\texttt{tests/test\_confound\_recovery.py}), and enforced by
the central gate orchestrator (\texttt{src/gates/orchestrator.py}) plus the
\texttt{Makefile} targets \texttt{test-portable} / \texttt{gate-e2e-synthetic}
/ \texttt{gate-orchestrate} that write immutable artifacts
(\texttt{artifacts/gates\_run\_*.json}) and support synthetic CI/repro. The
engine remains conceded plumbing, with a frozen-backbone successor path
documented separately. The standing engineering invariants are preserved:
G0-first manifest power (\texttt{power.json} emitted by Gate~0), early dedup and
vet filtering in Gate~1 before the cat-disjoint split (leak resistance), and the
MPS compute-correctness gate on the conceded engine. See the
\texttt{README} quickstart (including the \texttt{src.protocols} standalone
quickstart), \texttt{FINAL\_DIRECTION.md}, and the evaluation chapter for the
full gate-enforced contract.

The confound protocol, the reliability check, and the wrapper are shipped on an
explicitly \emph{conceded} engine: a frozen
DINOv2 ViT-S/14 backbone~\cite{oquab2023dinov2} (the register variant,
dinov2\_vits14\_reg, is the field default---registers suppress attention
artifacts that hurt the dense, localized per-AU features, A/B'd against the
plain variant before locking ViT-S/14) feeding five per-AU
rank-consistent CORN ordinal heads~\cite{shi2023corn}. We claim no novelty for
this engine and state so plainly throughout; it is plumbing that instantiates
the methods, not a contribution. The shippable spine is a
\emph{binary-plus-abstention} floor: a calibrated pain decision taken at a
welfare-asymmetric operating point fixed to pain-recall~$\geq 0.90$ (rather than
a symmetric-cost Youden-J or $F_1$ knee), accompanied by a defer-to-vet
abstention curve reported as a one-sided $95\%$ negative-predictive-value lower
bound (the word ``guaranteed'' is avoided as unsupportable at the budgeted
sample size). The wrapper's primary reported artifact --- cited supporting plumbing,
not a headline contribution --- is a harm-ratio-weighted decision curve that
sweeps the undertreat-to-overtreat asymmetry as a range. A strict circularity firewall holds throughout: the
sensitivity and specificity of the $0.39$ decision are estimated only on
vet-confirmed labels, and the kappa-versus-VLM agreement is never treated as
validation of that decision.

The proposal is governed by a self-honesty constraint that the abstract is
written to satisfy: every claim above must continue to hold if the exploratory
$0$--$10$ ordinal layer is dropped entirely. That ordinal layer is built and
inspected but is \emph{never} entered as a validated claim, and no validated
sentence in this work rests upon it. What we report in the results chapters is
therefore not a table of our accuracies but the evaluation \emph{protocol}, the
gate-conditioned expected outcomes, and the pre-registered decision rules
(power-calc budget, kappa CI-lower-bound floors, severity-cell collapse, and
the binary-plus-abstention floor as the likely v1 ship) under which each artifact ships or is demoted. This
document is therefore framed as a methods / protocol paper, demonstrated on
synthetic data and planted controls and targeted at a machine-learning
evaluation / trustworthiness or clinical-ML-methods venue rather than an
animal-welfare or veterinary journal as its first outlet; the empty data
directory, the missing independent vet anchor, the Roboflow-binary-only
substrate, and the non-reportable $n_{\text{pos}}=16$ operating point are
disclosed as limitations, with real-cat application named as future work. The
intended product is a decision-support triage instrument --- ``grimace
consistent with pain; recommend veterinary assessment'' --- and never an
autonomous analgesia trigger.

\vspace{1em}
\noindent\textbf{Keywords:} feline pain assessment; Feline Grimace Scale;
facial action units; vision--language model weak labeling; quadratic-weighted
kappa; trustworthy machine learning; calibration and abstention; decision-curve
analysis; confound attribution; ordinal regression.
